\documentclass[10pt,twocolumn]{article}
\usepackage[utf8]{inputenc}
\usepackage{amsmath,amssymb,amsfonts}
\usepackage{graphicx}
\usepackage{booktabs}
\usepackage{hyperref}
\usepackage{microtype}
\usepackage{geometry}
\geometry{margin=0.75in}

\title{\textbf{RAGeval: Claim-Level Evidence Support for Retrieval-Augmented Generation (Draft)}}
% Numbers in this file match reports/phase7/v1_matrix/metrics.json
% from run 20261028T183042Z\_a97ad68f. This is a draft. n=24 is underpowered.
% Do not cite as a confirmatory result. This project is not Niu et al. 2024 RAGTruth.

\author{
  \textbf{Aarav Khanal} \\
  Department of Computer Science
}
\date{\today}

\begin{document}
\maketitle

\begin{abstract}
RAGeval measures whether RAG answers are supported by retrieved evidence at the claim level. It is a software-and-pilot-data project, not a claim of novelty relative to FActScore, ALCE, or Niu et al.\ (2024) RAGTruth. The historical v1.0.0 mock-generator pilot (16 documents, 24 items) produced identical mixed-definition Recall@5 $=0.8333$ across six retrieval settings. A new matched real-generator companion on v1.3.0 (29 documents, 40 items) yields answerable-only Recall@5 $=1.0$ across the same six settings; claim-support rates range from $0.5917$ to $0.6500$. These different recall definitions cannot be compared directly. Items share documents, so differences are diagnostic. USD cost remains unreported without a provider rate; token usage is recorded. No overall score is used.
\end{abstract}

\section{Introduction}
Retrieval-augmented generation (RAG) can still emit unsupported or contradicted claims while looking fluent. RAGeval decomposes answers into claims, verifies them against retrieved chunks with DeBERTa-v3 NLI, and reports named metrics (claim support, contradiction, Recall@$k$, answer EM/F1, abstention). Faithfulness here is \emph{support by supplied evidence}, not world truth.

\section{Method (brief)}
Official generation in the recorded run is a blinded mock that must not read \texttt{test\_type}. Chunking is fixed-size MiniLM or BM25 or hybrid RRF, optional MiniLM cross-encoder rerank. Citation existence precision is separate from citation support precision. Mixed-support items keep a fully supported gold answer and an example mixed string for RQ3.

\section{Pilot results}
\label{sec:results}
All numbers in Table~\ref{tab:ablation} are from run \texttt{20261028T183042Z\_a97ad68f} (\texttt{reports/phase7/v1\_matrix/metrics.json}). Generator: mock. Verifier: NLI. Split: all 24 items.

\begin{table*}[t]
\centering
\small
\caption{Recorded v1.0.0 ablation. Recall@5 and MRR do not change across configs. Draft; $n=24$.}
\label{tab:ablation}
\begin{tabular}{lccccccr}
\toprule
\textbf{Configuration} & \textbf{Support} & \textbf{Unsupp.} & \textbf{Contra.} & \textbf{Recall@5} & \textbf{MRR} & \textbf{Abst.} & \textbf{Latency (ms)} \\
\midrule
Dense\_Chunk128 & 0.4375 & 0.1806 & 0.3819 & 0.8333 & 0.8125 & 0.000 & 2572 \\
Dense\_Chunk256 (Baseline) & 0.4375 & 0.1806 & 0.3819 & 0.8333 & 0.8125 & 0.000 & 2349 \\
Dense\_Chunk512 & 0.4375 & 0.1806 & 0.3819 & 0.8333 & 0.8125 & 0.000 & 2858 \\
Sparse\_BM25\_Chunk256 & 0.5000 & 0.2361 & 0.2639 & 0.8333 & 0.8125 & 0.000 & 2472 \\
Hybrid\_RRF\_Chunk256 & 0.5625 & 0.1806 & 0.2569 & 0.8333 & 0.8125 & 0.000 & 2192 \\
Hybrid\_RRF\_Reranked & 0.5833 & 0.3750 & 0.0417 & 0.8333 & 0.8125 & 0.000 & 1590 \\
\bottomrule
\end{tabular}
\end{table*}

\subsection{Corpus}
Short U.S.\ government public-domain excerpts (NASA Apollo 11, CDC influenza types, NPS Yellowstone) plus original CC-BY-4.0 fiction (HarborNet, LumenGrid, Nerexin, CedarRail, Millford). This is not SEC, NEJM, or OSDI source text, and not the Niu et al.\ RAGTruth corpus.

\subsection{Findings (inconclusive where $n$ is small)}
\begin{itemize}
    \item \textbf{RQ1.} Hybrid RRF did \emph{not} improve Recall@5 vs dense MiniLM on this 16-document corpus (relative lift $0\%$). Retrieval differences are not identifiable.
    \item \textbf{RQ2.} On 24 gold claim--premise pairs (\texttt{reports/verifier\_eval\_v1.md}), DeBERTa macro-F1 $0.6921$ exceeded a lexical heuristic $0.5140$. Partial-support F1 was $0.0$. Underpowered.
    \item \textbf{RQ3.} Mixed-support slice ($n=4$): token F1 $0.6292$ on every config; claim support $0.375$--$0.625$.
    \item \textbf{RQ4.} Abstention accuracy $0.0$. The blinded mock copies retrieved text when the index is non-empty.
    \item \textbf{RQ5.} Reranked mean latency was lower than dense-256 in this snapshot ($1590$ vs $2349$~ms). That does not confirm a latency ranking (caching). Cost is null.
\end{itemize}

Paired tests vs dense-256 are stored in \texttt{reports/phase7/v1\_matrix/significance.json}. All have \texttt{underpowered=true}. A raw $p=0.0109$ on hybrid vs baseline claim support is not treated as significant.

\section{Limitations}
Mock generator on the published table; $n=24$; current licensed corpus is v1.4.0 ($n=47$, 36 documents). The full set still shares documents. An independent companion slice now has $n=30$ items with disjoint gold documents; labels remain single-author, so that slice is not a confirmatory ranking. Single human annotator (a disclosed LLM rater B yields diagnostic kappa $0.9423$, not confirmatory IAA). Verifier partial-support F1 on the 24-pair set remains $0.0$. Historical reports used a mixed answerable/unanswerable Recall@5 definition. New metric schema v2 excludes no-gold questions from Recall@5 and reports empty retrieval on those questions separately; old and new recall values are not directly comparable. There is no overall score. USD cost is unavailable for the NVIDIA NIM trial endpoint because no published rate was identified; token usage is reported instead. Live NVIDIA remains a companion; the official system under test is the blinded mock on run \texttt{20261028T183042Z\_a97ad68f}. This file is a draft.

\section{Matched real-generator companion}
\label{sec:live}
The recorded \texttt{reports/phase7/live\_matrix\_v1/metrics.json} combines two completed runs (five retrieval settings and one reranker setting) over the same \texttt{benchmark\_v1} 1.3.0 questions. All six use \texttt{nvidia/nemotron-3.5-bolt-e4b}, temperature zero, top-five retrieval before optional reranking, and DeBERTa NLI. The 128-, 256-, and 512-word dense treatments remain identical on this short-document corpus; the separate \texttt{chunk\_benchmark\_v1} report tests actual chunk changes.

\begin{table*}[t]
\centering
\small
\caption{Diagnostic real-generator companion, 40 items sharing 29 documents. Recall@5 is answerable-only ($n=34$). USD price was unavailable; recorded token counts use API usage when returned and a local fallback otherwise.}
\label{tab:live}
\begin{tabular}{lccccrr}
\toprule
\textbf{Configuration} & \textbf{Support} & \textbf{Recall@5} & \textbf{MRR} & \textbf{Abst.} & \textbf{Input tok.} & \textbf{Output tok.} \\
\midrule
Dense 128 & 0.6375 & 1.0000 & 0.9265 & 1.000 & 19647 & 1372 \\
Dense 256 & 0.6375 & 1.0000 & 0.9265 & 1.000 & 19647 & 1372 \\
Dense 512 & 0.6375 & 1.0000 & 0.9265 & 1.000 & 19647 & 1371 \\
Sparse BM25 & 0.6500 & 1.0000 & 0.9618 & 1.000 & 19820 & 1358 \\
Hybrid RRF & 0.5917 & 1.0000 & 0.9216 & 1.000 & 19903 & 1468 \\
Hybrid+rerank & 0.6375 & 1.0000 & 0.9216 & 1.000 & 13555 & 1457 \\
\bottomrule
\end{tabular}
\end{table*}

All answerable items contain a gold document in the top five for every setting, so this small corpus cannot identify a Recall@5 improvement. Differences in MRR and claim support are descriptive. Six no-gold questions are excluded from Recall@5; empty retrieval occurred on one of them under each setting. The published v1.0.0 table above remains a historical mock-generator result with different metric semantics.

\section{Conclusion}
On this pilot, retrieval metrics did not move. Claim-level scores did move on mixed-support items while token F1 did not. Those observations are descriptive of run \texttt{20261028T183042Z\_a97ad68f}, not a ranking of RAG systems.

\bibliographystyle{plain}
\bibliography{references}

\end{document}
